\subsection{Vanishing of the Kummer shear on elliptic Galois slices}

\begin{claim}{FF.49}{Proved}\label{clm:FFellipticshear}
Let \(E/B_p\) be finite, let \(\varphi\) be a pinned
relative automorphism of \((\mathbb A^1_E)_{\acute et}\)
with identity actual geometric point function, and let
\[
                  \pi:C\longrightarrow\mathbb P^1_E
                                                        \tag{FF.49.1}
\]
be a finite Galois cover with smooth geometrically
integral genus-one source. On a branch-complement open,
use the unique actual scheme realization of
Claim~\ref{clm:FFlatteskernel} to regard its generic
slice action as an outer automorphism of \(G_{E(C)}\).
For every finite constant extension \(L/E\), every
\(m\ge1\), and every \(f\in L(C)^\times\), this
outer action satisfies
\[
                     \varphi_C^*\kappa_m(f)=\kappa_m(f).
                                                        \tag{FF.49.2}
\]
Equivalently it is the identity on the full abelian
profinite quotient of \(G_{L(C)}\), after every such
constant extension. This includes arbitrary translated
rational functions on the elliptic slice.

Thus every cyclic generic cover of this slice is
preserved as a covering object. Neither its actual
point-function realization nor identity on the entire
elliptic slice topos is asserted by this claim.
\end{claim}

\begin{proof}
We use integral Kummer characters throughout. Since
\(\mu_\infty\subset B_p\), the coefficient sheaves
\(\mathbb Z_\ell(1)\) are constant on every field
under discussion. The relative structural comparison
fixes these coefficients and all characters inflated
from the constant field. Compatible roots of a
nonzero rational function define its character
\(\kappa_{\ell^\infty}(f)\); reduction modulo
\(\ell^r\) is \(\kappa_{\ell^r}(f)\).
Inner changes in the generic action do not affect
these characters.

\emph{Completed-point compatibility on the elliptic slices.}
The geometric local theory used in
Claims~\ref{clm:FFgenerickernel} and
\ref{clm:FFgenuszerokernel} supplies a point
bijection on every smooth projective covering curve,
compatible with finite covering maps. For a
realized genus-one object we must check that this
bijection also agrees at the omitted boundary points
with the extension of its actual scheme realization.
Agreement on the branch complement alone does not
justify that assertion.

Use the dominating Latt\`es cover constructed in
Claim~\ref{clm:FFlatteskernel}. Thus, after finite
constants, the original genus-one Galois cover is a
quotient of
\(Z=B\to\mathbb P^1\) with deck group
\(B[N]\rtimes H\), where \(B\) is elliptic,
\(H\ne1\) consists of origin-preserving group
automorphisms, and \(B/H\) has genus zero. Choose
\(N\) divisible by the degree of the translation
quotient as in that construction, and increase it
so that
\[
              N^2>\sum_{1\ne h\in H}\deg(1-h).
                                                        \tag{FF.49.3}
\]
For every \(t\in B[N]\), put
\(H_t=\tau_tH\tau_{-t}\). The quotient
\(Z/H_t\) has genus zero. These quotients jointly
separate all geometric points of \(Z\). Indeed,
if two distinct points \(z,z'\) have the same
image in the quotient by \(H_t\), then
\[
                  z'=h(z)+(1-h)t
\]
for some \(h\in H\setminus\{1\}\). For each
fixed \(h\), the nonzero isogeny \(1-h\) has
at most \(\deg(1-h)\) solutions \(t\) to this
equation, even before restricting to \(B[N]\).
The union of these exceptional sets has size less
than \(N^2\), proving separation.

Claim~\ref{clm:FFgenuszerokernel} matches the
completed generic point bijection with the scheme
isomorphism on each \(Z/H_t\). On the branch
complement, naturality in
Claim~\ref{clm:FFlatteskernel} identifies the
quotient arrows with the scheme realization on
\(Z\); these scheme diagrams extend to their
smooth projective completions. The generic point
bijections are also compatible with every quotient
arrow. Joint separation therefore makes the
completed point bijection on \(Z\) equal to the
completion of its scheme realization. The finite
surjection from \(Z\) to the original genus-one
cover then gives the same equality there, by
surjectivity on geometric points. The construction
works for every genus-one Galois slice used below.
No assertion about separation of branches by the
projections of a normalized fibre product is used.

\emph{The possible error is geometrically unramified.}
For the moment work geometrically and choose an
origin on \(\bar C\), writing it as an elliptic
curve \(A\). Each finite collection of points,
functions and maps used below is defined after a
finite extension of constants, where the actual
scheme realizations in
Claim~\ref{clm:FFlatteskernel} are available and
natural. Their geometric cohomology groups are
unchanged by these extensions. The preceding
argument identifies their completed geometric
point functions as well.

Every geometric inertia exponent of the generic
base action is one, by
Claim~\ref{clm:FFallprimekernel}. The inertia at a
point of a finite covering curve is a finite-index
subgroup of the corresponding base inertia. After
the geometric point identification supplied by the
actual realization, the restricted exponent is
therefore also one. This holds at every point of
the smooth completion, including points over
infinity. Conjugator choices do not change the
conclusion: geometric curve inertia is procyclic,
and conjugation by its decomposition group acts
through the cyclotomic character, which is trivial
on the constant field here.

It follows that the geometric difference
\[
 b_1(f):=\varphi_C^*\kappa_{\ell^\infty}(f)
                         -\kappa_{\ell^\infty}(f)
       \ \in H^1(A,\mathbb Z_\ell(1))
                                                        \tag{FF.49.4}
\]
is unramified at all points. The indicated group
means geometric cohomology of the proper curve.
Indeed the residue of a Kummer character on point
inertia is the valuation of its function, and all
these residues of the difference vanish. The
resulting unramified character factors through the
proper curve fundamental group. It is an integral
class, not just a rational cohomology class.

\emph{Multiplication covers and their translations.}
For every positive integer \(n\), the map
\(\pi\circ[n]:A\to\mathbb P^1\), after finite
constants, is again a Galois cover with genus-one
source. To check this without assuming that all
deck transformations fix the origin, write each
deck transformation of \(\pi\) as
\(P\mapsto h(P)+t\), where \(h\) is a group
automorphism. It has exactly \(n^2\) lifts through
\([n]\), namely
\(P\mapsto h(P)+s\) with \(ns=t\). These give
\(n^2\deg(\pi)\) distinct automorphisms of the
composite, equal to its degree. All are defined over
one finite constant extension. In particular the
deck transformations include all translations
\(\tau_t\), \(t\in A[n]\).

Denote this covering object by \(A_n\), retaining
its map to the original rational base, and denote
its induced generic action by \(\varphi_n^*\).
Claim~\ref{clm:FFlatteskernel} realizes \(A_n\),
its deck arrows, and the map
\([n]:A_n\to A_1\) naturally on their branch
complements. Hence on cohomology
\[
 \varphi_n^*[n]^*=[n]^*\varphi_1^*,\qquad
 \varphi_n^*\tau_t^*=\tau_t^*\varphi_n^*.
                                                        \tag{FF.49.5}
\]
These equations are not an assumed extension to
unrelated rational subfields: all their arrows are
over the original base. The possible inner
ambiguities in the generic group comparisons are
irrelevant on \(H^1\).
Define \(b_n(h)\) as in (FF.49.4). The inertia
argument makes every such class lie in
\(H^1(A,\mathbb Z_\ell(1))\).
Translations act as the identity on this group,
and \([n]^*\) acts as multiplication by \(n\).
One may see both statements from
\(A(\mathbb C)=\mathbb C/\Lambda\): translations
are homotopic to the identity and multiplication
acts by \(n\) on the lattice. The algebraic and
profinite identifications follow from the
characteristic-zero comparison theorem
\citep[Expose~XII, Theoreme~5.1]{FF_SGA1}.
The group is free of rank two over \(\mathbb Z_\ell\).

\emph{An integral averaging identity.}
Suppose first that the principal divisor
\[
                       D=\operatorname{div}(f)
                              =\sum_i m_i[P_i]
\]
has some coefficient equal to \(1\) or \(-1\).
Choose \(Q_i\) with \(nQ_i=P_i\). Since \(D\)
is principal, \(\sum_i m_i=0\) and
\(\sum_i m_iP_i=0\) in the elliptic group.
Consequently \(s=\sum_i m_iQ_i\) belongs to
\(A[n]\). Adjust the chosen \(Q_i\) at a
coefficient \(\pm1\) by the appropriate element
of \(A[n]\) so that this sum becomes zero.
The divisor
\[
                         D_0=\sum_i m_i[Q_i]
\]
now has degree zero and zero sum in the elliptic
group, so it is principal, say
\(D_0=\operatorname{div}(h)\).
All these choices use only finitely many algebraic
constants. Divisors then give
\[
          f\circ[n]=c\prod_{t\in A[n]}\tau_t^*h
                  \qquad(c\in\overline{\mathbb Q}^{\times}).
                                                        \tag{FF.49.6}
\]
A constant has zero geometric Kummer character.
Apply the difference operator on \(A_n\) to
(FF.49.6), and use (FF.49.5). It gives the equality
in integral geometric cohomology
\[
           n\,b_1(f)
              =[n]^*b_1(f)
              =\sum_{t\in A[n]}\tau_t^*b_n(h)
              =n^2b_n(h).
                                                        \tag{FF.49.7}
\]
Torsion-freeness permits cancellation of \(n\).
Taking \(n=\ell^r\) for every \(r\ge1\) shows
that \(b_1(f)\) belongs to
\(\ell^r H^1(A,\mathbb Z_\ell(1))\) for every
\(r\), and hence is zero.

This proves the assertion for every rational
function, not just the ones with a coefficient
\(\pm1\). Given \(f\), choose a function \(q\)
with a simple zero outside the support of
\(\operatorname{div}(f)\), and with no cancellation
there. For example, take a degree-two function on
the elliptic curve and subtract a constant avoiding
its finitely many branch values and the images of
that support. Both \(q\) and \(fq\) have such
a simple zero. The proved case and additivity give
\(b_1(f)=b_1(fq)-b_1(q)=0\).
Thus the geometric restriction of every error in
(FF.49.4) vanishes, for every \(\ell\).

\emph{Removing the arithmetic error by the norm.}
Return to an arbitrary finite \(L/E\) and
\(f\in L(C)^\times\). Put
\(F=L(T)\), \(K=L(C)\), and
\(d=[K:F]\). The full integral difference
\[
       \delta=\varphi_C^*\kappa_{\ell^\infty}(f)
                         -\kappa_{\ell^\infty}(f)
          \in\operatorname{Hom}_{\mathrm{cts}}
                     (G_K,\mathbb Z_\ell(1))
                                                        \tag{FF.49.8}
\]
has zero restriction to the geometric kernel by
the preceding argument. It therefore factors as
an inflated character
\(\gamma:G_L\to\mathbb Z_\ell(1)\).
This factorization is at the full compatible
integral level; no conclusion is drawn merely
from separate finite-level classes.

Corestriction commutes with the induced generic
automorphisms, because the actual slice realization
identifies the covering map with \(\pi\). The
Kummer norm identity and
Claim~\ref{clm:FFallprimekernel} on the rational
base imply
\[
 \operatorname{cor}_{K/F}(\delta)
       =\varphi^*\kappa_{\ell^\infty}(N_{K/F}f)
                   -\kappa_{\ell^\infty}(N_{K/F}f)=0.
                                                        \tag{FF.49.9}
\]
Since \(\delta\) is the restriction of the
character inflated from \(\gamma\), the
restriction--corestriction identity identifies the
left side with \(d\gamma\), inflated to \(G_F\).
Inflation is injective, and
\(\operatorname{Hom}_{\mathrm{cts}}
(G_L,\mathbb Z_\ell(1))\) is torsion-free.
Therefore \(\gamma=0\), and \(\delta=0\).
This norm argument needs no rational point on
\(C_L\), so it also handles genus-one torsors
before choosing an origin.

Reducing the equality modulo every \(\ell^r\)
and taking the product over primes proves
(FF.49.2). Every finite Kummer character comes
from some \(f\) by Kummer theory, and
\(\mu_m\subset L\), so these equalities are
exactly identity on the full abelian profinite
quotient. This conclusion determines cyclic
covering objects; it does not supply the further
actual point comparisons required to identify
the nonabelian generic action or the whole slice
topos.
\end{proof}
